\rok{Новембар 2023.}{I}

Први практични тест из Алгоритама и структура података

\zadatak{1}{Проналажење елемента у реду (Queue Peek)}
Написати методу која имплементира стандардан начин за проналажење елемента са врха реда.

\textbf{Додатне напомене за решавање задатка:}
\begin{itemize}
    \item Алгоритам написати у оквиру закоментарисане методе \texttt{peek(int value)}.
    \item Поменута метода се налази у оквиру класе \texttt{Queue} која се налази у придруженом шаблону.
    \item Обезбедити код у \texttt{Main} методи за тестирање алгоритма.
\end{itemize}



\zadatak{2}{Померање 0 на крај двоструко спрегнуте листе}
Креирати функцију која ће проћи кроз двоструко спрегнуту листу и померити све 0 на крај. Алгоритам је неопходно написати са линеарном временском комплексношћу $O(n)$.

\textbf{Додатни захтеви за решавање задатка:}
\begin{itemize}
    \item Алгоритам мора да резултује у промени вредности постојеће двоструко спрегнуте листе.
    \item Могу се мењати само вредности, а могу се превезивати и комплетни чворови.
    \item Захтеван потпис методе:
    \begin{Verbatim}[fontsize=\footnotesize, numbers=left, xleftmargin=10mm]
public void MoveZerosToEnd()
    \end{Verbatim}
\end{itemize}

\textbf{Примери:}
\begin{itemize}
    \item \textbf{Улаз:} $9 \leftrightarrow -4 \leftrightarrow 0 \leftrightarrow 1 \leftrightarrow 0 \leftrightarrow 0 \leftrightarrow 7 \leftrightarrow 0$ \\
          \textbf{Излаз:} $9 \leftrightarrow -4 \leftrightarrow 1 \leftrightarrow 7 \leftrightarrow 0 \leftrightarrow 0 \leftrightarrow 0 \leftrightarrow 0$
    \item \textbf{Улаз:} $1 \leftrightarrow 0 \leftrightarrow 0 \leftrightarrow 1 \leftrightarrow 0 \leftrightarrow 0$ \\
          \textbf{Излаз:} $1 \leftrightarrow 1 \leftrightarrow 0 \leftrightarrow 0 \leftrightarrow 0 \leftrightarrow 0$
\end{itemize}



\zadatak{3}{Excessive број}
Креирање функције која проверава да ли је неки број \textit{excessive number}. Excessive број је број код којег сума свих његових делилаца не сме бити већа од самог броја. Алгоритам је неопходно написати са линеарном временском комплексношћу $O(n)$.

\textbf{Додатни захтеви за решавање задатка:}
\begin{itemize}
    \item Алгоритам треба да врати \texttt{bool} вредност.
    \item Цифре прослеђеног броја неопходно је сместити у низ. Уколико се користи динамички низ, користити функцију \texttt{ToArray()} да би се добио статички низ.
    \item Захтеван потпис методе:
    \begin{Verbatim}[fontsize=\footnotesize, numbers=left, xleftmargin=10mm]
public bool IsNumberExcessive(long number)
    \end{Verbatim}
\end{itemize}

\textbf{Примери:}
\begin{itemize}
    \item \textbf{Улаз 1:} 12 \\
          Објашњење: $1 + 2 + 3 + 4 + 6 = 16$ \\
          \textbf{Излаз 1:} \texttt{false}
    \item \textbf{Улаз 2:} 15 \\
          Објашњење: $1 + 3 + 5 = 9$ \\
          \textbf{Излаз 2:} \texttt{true}
\end{itemize}

\sablon{AISP2023_Test1_GrupaI}{2023/Novembar/GrupaI/AISP2023_Test1_GrupaI.linq}
